% ============================================================
%  SECCIÓN 1 – PLANTEAMIENTO DEL PROBLEMA
% ============================================================
\section{Planteamiento del Problema}

\subsection{Contexto}

La donación de sangre constituye un pilar fundamental de los sistemas de salud a escala mundial. De acuerdo con la Organización Mundial de la Salud (OMS), se requieren al menos entre 10 y 20 donaciones por cada 1\,000 habitantes para garantizar un suministro adecuado [2]; sin embargo, numerosos países de América Latina [3], incluido Colombia, se encuentran significativamente por debajo de ese umbral. En el departamento del Cauca, y particularmente en Popayán, la situación se agrava por la dispersión geográfica y por la ausencia de mecanismos ágiles que conecten a potenciales donantes con quienes requieren transfusiones de forma urgente.

\subsection{Problemática Identificada}

A partir del análisis del entorno se identificaron los siguientes problemas centrales:

\begin{enumerate}[label=\textbf{\arabic*.}, leftmargin=1.5cm]

    \item \textbf{Información dispersa y desactualizada.}
    Las convocatorias de donación circulan por redes sociales, volantes o comunicados institucionales sin un canal unificado. El ciudadano desconoce cuál fuente es oficial, cuándo caduca la necesidad o cuáles centros de salud están habilitados.

    \item \textbf{Desconexión entre la oferta y la demanda.}
    Cuando un paciente necesita una transfusión, su familia recurre a cadenas de mensajes informales (WhatsApp, Facebook) sin posibilidad de filtrar por compatibilidad de grupo sanguíneo ni por cercanía geográfica, lo que dilata los tiempos de respuesta.

    \item \textbf{Desinformación y mitos.}
    Creencias populares ---como que donar sangre debilita al organismo, afecta el peso corporal o compromete la vida sexual--- disuaden a potenciales donantes que, en condiciones de información correcta, estarían dispuestos a participar.

    \item \textbf{Falta de seguimiento y motivación.}
    El donante eventual no recibe retroalimentación sobre el impacto de su donación, carece de un historial accesible y no dispone de recordatorios sobre cuándo puede volver a donar, lo que reduce la tasa de donación recurrente.

    \item \textbf{Ausencia de herramientas para profesionales de salud.}
    Los centros de salud y hospitales no cuentan con un medio digital eficiente para publicar campañas de donación, gestionar cupos ni comunicar requisitos de manera estandarizada a la comunidad.

    \item \textbf{Comunicación insegura y no moderada.}
    Los canales informales de comunicación entre solicitantes y donantes no garantizan la privacidad de los datos ni limitan las interacciones a fines estrictamente relacionados con la donación.

\end{enumerate}

\subsection{Justificación}

Frente a la problemática descrita, surge la necesidad de una herramienta tecnológica que centralice el proceso de donación de sangre, conecte de manera inteligente a donantes y solicitantes con base en compatibilidad sanguínea y proximidad geográfica, y proporcione un entorno seguro y moderado de comunicación. Una solución móvil ---dado que el teléfono inteligente es el dispositivo de uso predominante en la población objetivo--- puede reducir las barreras de acceso, incrementar la tasa de donación voluntaria y, como consecuencia directa, contribuir a salvar vidas.
